\section{Conclusion}
Having engaged in the practical optimisation of the monomeric composition 
 of a simple bio-based branching thermoset polyester (poly (glycerol
 citrate itaconate)) with respect to an objective of interest
 (percentage mass loss via curing) we have discovered 
 an experimental global
 optimum of \(\Delta m_{\%}^{\circ}=20.77\%\) at the sample coordinates
 (\(s_{1}=0.337\),\(s_{2}=0.648\),\(b_{1}=0.773\)) corresponding to a
 stoichiometric balance of 1 glycerol to 0.77 citric acid to 0.37 itaconic
 acid.
This was discovered by a sequential Bayesian technique whose
 sampling policy adopted the well known probability of improvement
 acquisition function.

The utility of Bayesian techniques in sample-efficient global optimisation
 was proven through 10,000 simulations of their behaviour across five
 different surrogate functions emulating the objective function.
In the case of emulator GP-Mat-05 the expected improvement and probability
 of improvement sequential techniques were found to exhibit a
 65.0\% and 23.8\% probability of uncovering the computational 
 `region of global optimality', whist both the pseudorandom, quasirandom,
 and objective-naive sequential techniques displayed much lower 
 probabilities (0.6\%, 1.3\%, and 0.5\% respectively).

If the intention is for poly(glycerol citrate itaconate) to perform 
 the role of a useful thermoset resin for moulding applications, 
 future work would have to focus on the mitigation of the bubbling 
 effect, improvement of mechanical strength, and reductions in 
 the cure time.
\textcolor{RawSienna}{
The last of these is of particular importance given the relatively
 higher amounts of process-energy required to prepolymerise and
 cure the presently described branched thermosets when compared
 with conventional unsaturated polyester resins (UPR).
Modelling of these two processes with error propagation as 
 described in the supporting information (see 
 Section~\ref{subsec:NoPEC}) suggested that conventional UPR
 thermosets require 81\% less process energy than the 
 branched-polyester route, with a 
 5\textsuperscript{th}–95\textsuperscript{th} percentile 
 uncertainty range of 57–91\%.
This difference was found to be dominated by the prolonged 
 90~\unit{\degreeCelsius} oven residence time required for the 
 branched polyester.
The estimates remained sensitive to oven power, batch loading, cure
 duration and the process boundaries applied to prepolymer
 manufacture.
}

\textcolor{RawSienna}{We therefore suggest that future work focuses upon 
 dual-objective optimisation problems which consider both mechanical 
 strength maximisation alongside oven residence time minimisation, 
 whilst adopting predictor variables such as prepolymer stoichiometry,
 temperature and pressure, and quantities of catalytic reagent
 added.
}
Alternatively, different bio-based and readily-available monomers 
 could be considered (e.g. succinic acid) which may encourage 
 development of novel thermoset networks that might impact upon the 
 objectives of interest.
Note that addition of further monomers to the stoichiometric balancing
 mechanism described in this paper would, due to the curse of
 dimensionality, require many more samples to effectively explore and
 optimise these systems: this would encourage the replacement of 
 human-in-the-loop experimentation with laboratory automation.

\sout{
Furthermore, this sequential Bayesian technique and a second (which employed
 the alternative expected improvement acquisition function) displayed
 better sample-efficiency than simple one-shot techniques in the field.

Utility of these Bayesian techniques in sample-efficient global
 optimisation was proven through 1,000 simulations of their behaviour
 when sampling a surrogate function emulating the objective function; the
 `expected improvement' and 'probability of improvement' sequential
 techniques were found to exhibit a 5.3\% and 1.1\% probability of
 uncovering the experimental region of global optimality, whilst both the
 pseudorandom and quasirandom one-shot techniques entirely failed to
 sample these regions.
}